\section{对称性与对应}\label{symq}

\subsection{概览}\label{oversym}

如第 \ref{problem} 节所解释，对于带预量子线丛 $\frL$ 的实辛流形 $M$，不可能以这样的方式量子化：使经典力学的自然对称群 $\G$ 实现为量子 Hilbert 空间 $\H$ 上的酉变换群。（回忆 $\G$ 的元素是辛同胚 $\varphi:M\to M$ 连同 $\varphi$ 到 $\frL$ 的对称的提升。）既然如此，我们不如问如下问题：对于给定的量子化方法，$\G$ 的哪个子群在量子力学上实现为酉对称群？

在几何量子化中，这个问题有标准答案。$M$ 的几何量子化依赖于 $M$ 的极化 $\P$ 的选择。对给定的 $\P$，在量子力学上实现为对称群的 $\G$ 的子群，是保持 $\P$ 的子群 $\G_\P$。在有时重要的特殊情形中，几何量子化定义的 Hilbert 空间 $\H$ 容许 $\G$ 的更大子群以酉方式作用；但几何量子化的一般答案是，在量子理论中起作用的经典对称子群是 $\G_\P$。

例如，若 $M=T^*N$ 是带标准辛结构的余切丛，$\P$ 是叶为余切丛纤维的实极化，则我们可以相当显式地描述 $\G_\P$。保持 $\P$ 不变的辛同胚群 $\G_\P$ 由两类辛同胚生成。第一，$N$ 的每个微分同胚都延拓为 $T^*N$ 的保持辛结构的微分同胚。这些微分同胚置换 $\P$ 的叶。此外，还有把每个叶映到自身的 $T^*N$ 的辛同胚。若 $q^i$ 是 $N$ 上的局部坐标，$p_i$ 是余切丛纤维上相应的线性函数，则把每个叶映到自身的辛同胚保持 $q^i$ 不变，并把 $p_j$ 平移为
$p_j\to p_j+\partial f/\partial q^j$，
其中 $f(q^1,q^2,\cdots,q^n)$ 是某个函数（称为正则变换的生成函数）。

膜量子化中的对称群是什么？在膜量子化中，第一步是把带辛形式 $\omega$ 的实辛流形 $M$ 复化为带全纯辛形式 $\Omega$ 的复辛流形 $Y$。$Y$ 必须使得：若把 $Y$ 视为带辛形式 $\Im\Omega$ 的实辛流形，则其 $A$-模型良定义。（例如，$Y$ 可以是完备超 Kähler 流形，如第 \ref{amo} 节所讨论。）选定合适的 $Y$ 后，定义带 $\CP$ 丛 $\L$ 的拉格朗日膜 $\B$，则膜量子化的量子 Hilbert 空间为 $\H=\Hom(\B,\B_\cc)$。

对膜量子化的对称群，有一个显然的候选答案。$Y$ 的 $A$-模型的自然对称，是保持辛形式 $\Omega$ 的 $Y$ 到自身的可逆全纯映射，即全纯辛同胚。若 $Y$ 的全纯辛同胚把 $M$ 映到自身，且其在 $M$ 上的作用提升为对 $\frL$ 的作用，则它可能是膜 $\B$ 的对称。这给出膜量子化对称的自然候选。设 $\G_Y$ 为 $\G$ 中这样的元素构成的子群：其在 $M$ 上的作用可解析延拓为 $Y$ 的全纯辛同胚。则 $\G_Y$ 是酉作用于 $\H=\Hom(\B,\B_\cc)$ 的 $\G$ 的子群的自然候选。

尽管这个提法可能是正确的，要证明它实际上有技术困难。为使 $A$-模型具体化，通常要在 $Y$ 上选取 Riemann 度量 $g$。这可以是第 \ref{amo} 节所讨论的完备超 Kähler 度量。至少，$g$ 应与实辛形式 $\omega_K=\Im\Omega$ 一起定义一个近复结构 $K$，使 $\omega_K$ 为 $(1,1)$ 型且正。此外，$g$ 应使 $Y$ 的 $A$-模型真正存在。只要 $g$ 在使 $A$-模型存在的容许度量类中变化，$Y$ 的 $A$-模型就与 $g$ 无关。遗憾的是，最后这个条件很难处理，其全部含义也不易理解。

现在设 $\varphi$ 是 $Y$ 的全纯辛同胚。一般说来 $g$ 不是 $\varphi$-不变的，而会被 $\varphi$ 变成另一个度量 $g^\varphi$。带度量 $g^\varphi$ 的 $Y$ 的 $A$-模型一般可能并不良定义。即使良定义，要论证 $g$ 与 $g^\varphi$ 给出等价的 $Y$ 的 $A$-模型，我们也需要知道可以在容许度量类中把 $g$ 与 $g^\varphi$ 插值。总可以用一族与合适近复结构相伴的度量在 $g$ 与 $g^\varphi$ 之间插值。\footnote{其论证与稍后关于超 Kähler 度量的一个论证类似。$Y$ 的辛结构 $\omega_Y$ 把 $Y$ 切丛的结构群从 $\GL(4n,\R)$ 约化到 $\Sp(4n,\R)$。使 $\omega_Y$ 为正且为 $(1,1)$ 型的度量，对应于把结构群进一步约化到 $\Sp(4n,\R)$ 的极大紧子群 $\UU(2n)$。可以在任意两个这样的约化之间插值，因为商 $\Sp(4n,\R)/\UU(2n)$ 可缩。}但并不显然这种插值能在使 $A$-模型存在的度量类中进行。

物理上，人们更愿意把 $A$-模型看作紫外完备的 $\sigma$-模型的扭转版本。试图考虑基于诸如 $g^\varphi$ 的度量的 $A$-模型，一般会把我们带出那个世界；而基于从 $g$ 到 $g^\varphi$ 插值的论证可能把我们带得更远。

为进一步给读者定位，我们将描述一些具体问题，说明这个问题在何处成问题、在何处不成问题。然后我们将描述另一种途径，使人们能在紫外完备的世界中研究辛同胚。

我们将在 $M=\R^{2n}$ 的量子化背景下说明这个故事的两个方面，其中实值坐标为 $\vec x=(x^1,x^2,\cdots,x^{2n})$，标准辛结构为 $\omega=\frac12\sum_{i,j=1}^{2n}\omega_{ij}\d x^i\d x^j$。首先，我们试着用膜量子化解释标准结果\footnote{这也可由第 \ref{comparing} 节的论证推出：$\R^{2n}$ 的膜量子化不依赖于极化的选择，等价于用实线性极化的 $\R^{2n}$ 的几何量子化。}：$M$ 的量子化可以选为在仿射线性辛同胚群 $\vec x\to A\vec x+\vec b$ 下不变，如第 \ref{problem} 节所讨论。我们以熟悉的方式开始，把 $M$ 复化为复辛流形 $Y=\C^{2n}$，此时 $x^i$ 是 $Y$ 上的全纯坐标。$Y$ 的复结构与复辛结构当然是平移不变的，即在 $\vec x\to\vec x+\vec b$（$\vec b$ 为复常数）下不变。按第 \ref{amo} 节的意义，$Y$ 的平移不变复辛结构可通过选取合适的平移不变 Riemann 度量 $g$ 延拓为平移不变超 Kähler 结构。这样的 $g$ 构成的空间可如下描述。由于假设了平移不变性，我们可在一点 $y\in Y$ 处作分析。$Y$ 在 $y$ 处的切空间维数为 $4n$，故其线性变换群是 $\GL(4n,\R)$ 的一个拷贝。$Y$ 的复辛结构在子群 $\Sp(2n,\C)\subset\GL(4n,\R)$ 下不变。延拓其复辛结构的 $Y$ 上平移不变超 Kähler 度量，不在 $\Sp(2n,\C)$ 下不变，而在其极大紧子群下不变。这样的子群同构于辛群的紧实形式，我们称之为 $\Sp_c(2n)$。$\Sp(2n,\C)$ 的任一 $\Sp_c(2n)$ 子群都是被某个超 Kähler 度量保持的子群，且所有这样的子群在 $\Sp(2n,\C)$ 中共轭。因此，延拓 $Y$ 的复辛结构的平移不变超 Kähler 度量的空间是齐性空间 $\sW=\Sp(2n,\C)/\Sp_c(2n)$ 的一个拷贝。由 $\sW$ 参数化的任何超 Kähler 度量都可用于定义 $Y$ 的 $A$-模型。$\sW$ 可缩，特别地连通且单连通。故可在 $\sW$ 中任意两点 $w_1,w_2\in\sW$ 之间插值，而且这一插值在同伦意义下唯一。（唯一性很重要，因为从 $w_1$ 到 $w_2$ 的同伦不同的插值可能导致相差一个酉变换的量子理论之间的等价。）因此，若我们用平移不变超 Kähler 度量定义 $Y$ 的 $A$-模型，用哪一个并无关系。

现在我们可以用新的方式理解为什么 $M$ 的膜量子化在 $M$ 的仿射线性变换下不变。$M$ 的仿射线性辛同胚 $\vec x\to A\vec x+\vec b$ 可解析延拓为 $Y$ 的全纯辛同胚，由同一公式定义。一般的仿射线性辛同胚不会保持用于定义 $Y$ 的 $A$-模型的平移不变超 Kähler 度量。但它会把这个平移不变超 Kähler 度量变成另一个平移不变超 Kähler 度量，而 $Y$ 的 $A$-模型并不在乎我们用哪一个平移不变超 Kähler 度量。因此，通过 $A$-模型的 $M$ 的量子化在仿射线性辛同胚下不变；仿射线性辛同胚群作用在量子 Hilbert 空间 $\H$ 上。

为看出这类论证的局限，我们注意以下事实。$M=\R^{2n}$ 的、能延拓为 $Y$ 的全纯辛同胚的辛同胚群，比仿射线性辛同胚群大得多。在本讨论中，我们只考虑多项式辛同胚，因为我们在 $A$-模型中只允许多项式函数。为构造一个非常大的多项式辛同胚群，先把变量 $x^1,\dots,x^{2n}$ 分成 Poisson 交换的坐标 $q^1,\cdots,q^n$ 与共轭动量 $p_1,\cdots,p_n$。对任何多项式函数 $f(q^1,\cdots,q^n)$，考虑把 $p_i,q^j$ 映为
\begin{align}\label{worfo} p'_i&=p_i+\frac{\partial f}{\partial q^i},\cr q^i{}'&=q^i \end{align}
的辛同胚。这类辛同胚连同线性辛同胚\footnote{\label{enough}这里考虑线性辛同胚 $\vec x\to A\vec x$ 就足够了，不必考虑仿射线性辛同胚 $\vec x\to A\vec x+\vec b$，因为 $p$ 的常数平移已包含在式 \eqref{worfo} 中；而一旦包含线性辛同胚（它们可以交换 $p$ 与 $q$），$q$ 的常数平移也就包含在内了。}生成一个称为 $\R^{2n}$ 的驯顺辛同胚群的群。我们记之为 $\TAut(\R^{2n})$。为看出 $\TAut(\R^{2n})$ 有多么巨大，可以想象试着给这个群配一个 Lie 代数。（在这个无穷维情形中，这在技术上有问题。）这样的 Lie 代数将不得不把所有 $p$ 与 $q$ 的多项式函数作为生成元包含进来。确实，式 \eqref{worfo} 中的变换以 $f(q^1,\cdots,q^n)$ 为生成元。由于 $\TAut(\R^{2n})$ 包含线性辛同胚，而后者特别可以交换 $p$ 与 $q$，一个以任意多项式函数为生成元的假想 Lie 代数也必须包含 $p$ 的任意多项式函数。但这样一来，要得到在 Poisson 括号下封闭的函数类，我们就必须允许 $p$ 与 $q$ 的任意多项式。所以 $\TAut(\R^{2n})$ 大到这样的程度：若它能被当作 Lie 群处理，其 Lie 代数将包含正则变量的一切多项式。

对 $n=1$，已知 $\TAut(\R^{2n})$ 与 $\R^{2n}$ 的全体多项式辛同胚群 $\Aut(\R^{2n})$ 重合。对 $n>1$ 是否如此尚不清楚。（注意在 Lie 代数层面，人们看不出 $\TAut(\R^{2n})$ 与 $\Aut(\R^{2n})$ 的区别，因为如上所述，$\TAut(\R^{2n})$ 的 Lie 代数必须具有任意多项式生成元。）

在形变量子化中，我们可以考虑 $\C^{2n}$ 的一切全纯多项式辛同胚群 $\Aut(\C^{2n})$，以代替 $\Aut(\R^{2n})$。同样，我们可以考虑相应的驯顺子群 $\TAut(\C^{2n})$，以代替 $\TAut(\R^{2n})$；它由 $\vec x\to A\vec x+\vec b$ 连同式 \eqref{worfo} 中的辛同胚生成，只不过现在 $A$、$\vec b$ 与 $f$ 都可以是复的。这些群是 $\C^{2n}$ 上多项式函数经典环 $\A_0$ 及其 Poisson 括号的对称群。形变量子化把 $\A_0$ 换成非交换环 $\A$，有时称为 Weyl 代数（由生成元 $x_i$ 及关系 $[x_i,x_j]=\i\hbar\Omega_{ij}$ 定义）。经典对称群 $\TAut(\C^{2n})$ 与 $\Aut(\C^{2n})$ 的量子类比，是 $\A$ 的自同构群。$\TAut(\C^{2n})$ 与 $\Aut(\C^{2n})$ 同其量子类比有什么关系？

Kanel-Belov 与 Kontsevich \cite{BKK} 建立了惊人的结果：$\TAut(\C^{2n})$ 是 Weyl 代数自同构群的子群；他们并猜想，可能更大的群 $\Aut(\C^{2n})$ 就是 Weyl 代数的自同构群。这后来由 Kanel-Belov、Elishev 与 Yu \cite{BEY} 证明。（这些证明大量使用模素数 $p$ 约化；Weyl 代数在这样的约化后具有相当不寻常的性质。）正如第 \ref{symmeq} 节末尾将解释的，\cite{BKK,BEY} 的结果对 $\R^{2n}$ 的量子化（而非 $\C^{2n}$ 的形变量子化）有直接的类比。特别地，$\Aut(\R^{2n})$ 的中心扩张作为酉对称群作用在量子化 $\R^{2n}$ 所得的 Hilbert 空间上。这里没有量子形变。一方面，我们可能预期如此，因为这个中心扩张正是本讨论开头引入的群 $\G_Y$。另一方面，这个结果很引人注目，因为在 Lie 代数层面没有类似的东西成立。如前面所解释，$\Aut(\R^{2n})$ 甚至 $\TAut(\R^{2n})$ 的 Lie 代数将不得不包含 $\R^{2n}$ 上的一切多项式函数。由于从经典力学过渡到量子力学时通常的反常（第 \ref{problem} 节与附录 \ref{anomaly} 讨论过），这样的 Lie 代数无法在没有量子形变的情况下作用在量子 Hilbert 空间上。

用膜量子化恢复 \cite{BKK,BEY} 的部分结果是合理的目标。然而，$Y$ 的一般多项式辛同胚 $\varphi$ 会把 $Y$ 上的平移不变度量 $g$ 变成远离平移不变的度量 $g^\varphi$。一般说来，$g^\varphi$ 可能不给出紫外完备的 $\sigma$-模型；即使给出，我们也不知道如何在导致这类模型的 $Y$ 的度量类中于 $g$ 与 $g^\varphi$ 之间插值。因此，要从膜量子化恢复 \cite{BKK,BEY} 的部分结果，需要另一种途径。在第 \ref{corresp} 节中，我们将引入一种途径，使我们能在留在紫外完备世界的同时研究一般多项式辛同胚。这一途径的另一优点是，它为研究比辛同胚更一般的算符提供了良好的框架。

\subsection{界面与算符}\label{corresp}

辛流形 $Y$ 的辛同胚 $\varphi$ 是保持 $Y$ 的辛结构 $\omega$ 的可逆映射 $\varphi:Y\to Y$，即 $\varphi^*(\omega)=\omega$。具体地说，若 $x^i$ 是 $Y$ 上的局部坐标，则 $\varphi$ 可由把 $x'{}^{i}=\varphi^*(x^i)$ 表为 $x^k$ 的函数来定义：
